    % ==============================================================================
% BAB 1: PENDAHULUAN
% ==============================================================================

\chapter{Pendahuluan}
\label{chapter:pendahuluan}
\section{Latar Belakang}
Paradoks pangan global saat ini ditandai oleh kontradiksi tajam antara tingginya surplus produksi yang berujung pada pemborosan masif dengan krisis ketahanan pangan. Hal ini sejalan dengan urgensi global untuk mendaur ulang dan memanfaatkan kembali sisa pangan secara esensial demi menjaga keberlanjutan lingkungan \cite{raiEssentialRecyclingRepurposing2025a}. Di Indonesia, fenomena makanan yang berlebih kerap kali hanya terkumpul menjadi surplus, yang jika gagal disalurkan akan membusuk dan berubah status menjadi timbulan sampah makanan (\textit{food waste}). Permasalahan ini sangat masif, di mana 23 hingga 48 juta ton sampah makanan dihasilkan setiap tahunnya yang memicu kerugian ekonomi mencapai triliunan rupiah \cite{immawanFoodWasteIndonesia2024a}. Di sisi lain, inisiatif penyelamatan surplus makanan ini memiliki potensi besar untuk memperkuat perekonomian Usaha Mikro, Kecil, dan Menengah (UMKM) dengan memberikan solusi saling menguntungkan (\textit{win-win solution}) bagi produsen dan konsumen \cite{yulianaPotensiGerakanFood2022b}. Ironisnya, pemborosan ini terjadi berdampingan dengan tingginya ancaman kerawanan pangan dan prevalensi \textit{stunting} (gizi buruk kronis) yang masih menyentuh angka 32\% di berbagai wilayah agraris seperti di Kabupaten Bondowoso \cite{prayitnoExploringRoleFood2025b}. Krisis ini menuntut adanya intervensi strategi ketahanan pangan yang berkelanjutan, sebagaimana yang mendesak untuk diupayakan di berbagai wilayah seperti Jawa Tengah guna menangani masalah inflasi, kemiskinan, serta kekurangan gizi secara komprehensif \cite{pujiyantoStrategiKetahananPangan2025b}. Oleh karena itu, menyelamatkan surplus makanan pada tingkat konsumsi menjadi strategi mitigasi yang sangat krusial untuk menjaga kelestarian lingkungan dan menekan jejak karbon \cite{liuExploringHouseholdFood2023b}.

Meskipun inisiatif pembagian surplus makanan memiliki potensi penyelamatan yang besar, pelaksanaannya di lapangan masih terhambat oleh inefisiensi logistik dan kendala operasional. Makanan sisa yang gagal disalurkan ke tingkat rumah tangga terbukti memberikan dampak ekologis yang buruk, yang meliputi pencemaran lingkungan dan peningkatan emisi gas rumah kaca \cite{alliffiyaalkhansyaharsyDampakLimbahMakanan2025c}. Fakta ini divalidasi oleh analisis pada sektor rumah makan berskala besar di Kota Bogor, di mana tingginya kuantitas sisa makanan komersial berbanding lurus dengan besarnya pelepasan jejak emisi karbon \cite{fernandezCarbonFootprintAnalysis2024c}. Di sisi lain, kebijakan pengelolaan sampah sering kali terkendala oleh faktor perilaku, seperti kurangnya perencanaan berbelanja dan motivasi ekonomi yang menjadi pendorong utama \textit{food waste} di tingkat masyarakat \cite{lestariKebijakanPengelolaanSampah2022c}. Kesiapan ekosistem di Indonesia dalam melakukan pemanfaatan kembali (\textit{valorization}) sisa pangan pun masih terhambat oleh keterbatasan teknologi, tingginya biaya, dan kecemasan terhadap jaminan standar keamanan produk turunannya \cite{immawanFoodWasteIndonesia2024a}. Padahal, kampanye penyelamatan makanan dan edukasi kesadaran lingkungan terbukti sangat efektif sebagai faktor terpenting untuk memitigasi emisi jejak karbon, sebagaimana yang ditunjukkan secara empiris melalui studi pada kelompok institusi pendidikan tinggi di Tiongkok \cite{qianFoodWasteAssociated2022a}.

Untuk memotong rantai perubahan surplus menjadi sampah makanan sekaligus menambal celah inefisiensi logistik dan tingginya ancaman keamanan pangan tersebut, \textit{startup} JagoAI hadir menawarkan sebuah platform agregator logistik digital. Aplikasi ini dirancang secara khusus untuk menjembatani para penyumbang makanan berlebih dengan fasilitas sosial maupun masyarakat rentan secara \textit{real-time} melalui pemanfaatan Kecerdasan Buatan (AI). Dalam sistem pengenalan visualnya, pemanfaatan informasi kontekstual pada \textit{Computer Vision} terbukti sangat esensial untuk memperjelas dan menganalisis data visual sehingga mampu meningkatkan presisi serta efektivitas saat mendeteksi objek kelayakan pangan \cite{jamaliContextObjectDetection2025b}. Meskipun demikian, model AI pada sistem ini dirancang dengan sangat tangguh guna menghindari fenomena bias \textit{Clever Hans effect}, yaitu kegagalan kritis di mana model AI mencapai kinerja tinggi hanya dengan memanfaatkan korelasi palsu (\textit{spurious correlations}) beserta artefak pada latar belakang data, alih-alih mengandalkan hubungan kausal objeknya \cite{pathakUnmaskingCleverHans2026a}. Selain itu, arsitektur pendeteksi objek ini juga dimitigasi secara presisi agar tidak mudah dikecoh oleh metode manipulasi halus seperti \textit{adversarial patches} yang dapat menyebabkan misklasifikasi objek \cite{hopsonStudyAIObject2026b}. Dengan pondasi intelijensi yang tangguh, fitur AI pada platform ini menjadi instrumen otonom yang sangat ideal untuk memvalidasi kelayakan makanan secara mandiri sebelum disalurkan.

Guna memastikan seluruh proses distribusi berjalan aman dan transparan, platform ini dibekali dengan berbagai rincian fitur cerdas yang terintegrasi. Sistem ini menggunakan layanan \textit{Google Gemini Vision API} sebagai ``Polisi Pangan'' untuk memverifikasi keamanan, higienitas visual, dan kehalalan makanan donasi secara otomatis sebelum dipublikasikan kepada publik. Khusus bagi pengguna donatur korporat, tersedia pula fitur AI generatif eksklusif seperti \textit{Eco Packaging Editor} untuk merancang desain kemasan ramah lingkungan, serta \textit{CSR Copywriter} untuk otomatisasi pembuatan naskah kampanye sosial. Di lapangan, proses penjemputan makanan digerakkan oleh armada relawan kurir yang dimotivasi oleh sistem gamifikasi berbasis papan peringkat (\textit{leaderboard}). Seluruh proses serah terima ini diwajibkan menggunakan pemindaian kode QR secara nirkontak (\textit{contactless}) guna mencegah terjadinya klaim ganda. Akhirnya, setiap aktivitas donasi yang berhasil disalurkan akan direkapitulasi secara komprehensif ke dalam Dasbor Dampak Sosial (SROI), yang secara waktu nyata melacak kuantitas porsi makanan terselamatkan sekaligus mengalkulasi pencapaian pengurangan emisi jejak karbon bumi akibat pencegahan terbentuknya sampah makanan.


\section{Rumusan Masalah}
Berdasarkan pemaparan latar belakang di atas, maka rumusan masalah dalam pengerjaan Tugas Akhir ini difokuskan pada tiga poin pertanyaan berikut:
\begin{enumerate}
\item Bagaimana mendistribusikan surplus makanan dari donatur ke fasilitas sosial atau masyarakat rentan secara \textit{real-time} untuk mencegah makanan tersebut membusuk menjadi sampah (\textit{food waste})?
\item Bagaimana memvalidasi kelayakan dan keamanan surplus makanan secara visual dan otonom, serta memfasilitasi pembuatan konten kampanye keberlanjutan bagi donatur korporat?
\item Bagaimana memotivasi partisipasi relawan kurir, mencegah penyalahgunaan klaim ganda saat serah terima di lapangan, dan melacak pencapaian reduksi emisi karbon secara transparan?
\end{enumerate}

\section{Tujuan}
Berdasarkan rumusan masalah yang telah diuraikan sebelumnya, maka tujuan dari pengerjaan Tugas Akhir ini adalah sebagai berikut:
\begin{enumerate}
    \item Membangun platform agregator logistik digital berbasis \textit{Progressive Web App} (PWA) untuk mendistribusikan surplus makanan dari donatur ke fasilitas sosial atau masyarakat rentan secara \textit{real-time} guna mencegah makanan tersebut membusuk menjadi sampah (\textit{food waste}).
    \item Mengimplementasikan teknologi \textit{Computer Vision} (\textit{Google Gemini Vision API}) dan AI Generatif (\textit{Eco Packaging Editor} dan \textit{CSR Copywriter}) untuk memvalidasi kelayakan dan keamanan surplus makanan secara visual dan otonom, serta memfasilitasi pembuatan konten kampanye keberlanjutan bagi donatur korporat.
    \item Mengintegrasikan sistem gamifikasi (\textit{leaderboard}), pemindaian Kode QR nirkontak (\textit{contactless}), dan Dasbor Dampak Sosial (SROI) untuk memotivasi partisipasi relawan kurir, mencegah penyalahgunaan klaim ganda saat serah terima di lapangan, dan melacak pencapaian reduksi emisi karbon secara transparan.
\end{enumerate}

% \section{Cakupan Pengerjaan}
% Cakupan pengerjaan Tugas Akhir ini sebagai berikut:
% \begin{enumerate}
%     \item Ruang Lingkup Aplikasi: Sistem yang dibangun adalah aplikasi berbasis web (Progressive Web App/PWA) yang berfungsi khusus sebagai agregator logistik untuk memfasilitasi donasi makanan surplus halal yang masih layak konsumsi secara gratis atau non-komersial
%     \item Batasan Pengguna (Aktor): Sistem membatasi interaksi pada 6 (enam) peran pengguna yang saling terintegrasi, yaitu:
%     \begin{enumerate}
%         \item Donatur Individu: Masyarakat perorangan atau usaha mikro yang bertugas mengunggah stok makanan surplus skala kecil, serta wajib memindai Kode QR Penerima saat serah terima.
%         \item Donatur Kelompok/Korporat: Pelaku bisnis HORECA (skala besar) yang mengelola stok donasi dengan hak akses eksklusif terhadap fitur AI Generatif (pembuatan desain kemasan dan copywriting media sosial) untuk mendukung kampanye Corporate Social Responsibility (CSR).
%         \item Mitra Penerima: Masyarakat rentan atau yayasan yang bertugas mencari donasi, melakukan klaim (booking) untuk mengunci stok, dan menyajikan Kode QR miliknya untuk dipindai saat makanan tiba.
%         \item Relawan: Armada logistik sosial (crowdsourcing) yang menjemput dan mengantar makanan, serta memegang otoritas untuk wajib memindai (mandatory scan) Kode QR Penerima sebagai bukti sah penyelesaian distribusi.
%         \item Admin Pengelola: Pihak manajerial (seperti Dinas Sosial atau pengurus Food Bank Nasional) yang bertugas memverifikasi pendaftaran akun, menindaklanjuti laporan negatif pengguna, dan memantau Dasbor Analitik Dampak (SROI).
%         \item Super Admin: Tim pengembang (IT) yang memiliki akses backend tertinggi untuk melakukan pemeliharaan (maintenance) server dan menjaga stabilitas basis data.
%     \end{enumerate}
%     \item Batasan Fitur dan Teknologi Utama: Pengembangan sistem difokuskan pada manajemen distribusi digital yang telah dilengkapi dengan prosedur keamanan wajib pindai (mandatory scan) Kode QR saat serah terima logistik dan fitur gamifikasi Poin Sosial. Dari segi teknologi inti, aplikasi ini mengintegrasikan API Gemini 3     Engine yang mengeksekusi 4 (empat) perluasan fungsi kecerdasan buatan (AI), yaitu:
%     \begin{enumerate}
%         \item AI Quality Verification: Fitur pemrosesan citra (Computer Vision) yang bertindak sebagai bantuan skrining awal untuk memprediksi kelayakan fisik dan higienitas visual makanan.
%         \item AI Generatif (Eksklusif Korporat): Modul khusus bagi donatur tingkat perusahaan untuk mendukung kampanye Corporate Social Responsibility (CSR), yang mencakup fitur pemindaian bahan untuk rekomendasi resep, pembuatan (generate) gambar desain kemasan, serta penyusunan teks (copywriting) promosi media sosial.
%         \item AI Impact Analytics: Dasbor pelaporan otomatis yang merekam riwayat transaksi guna mengkalkulasi pencapaian dampak sosial (Social Return on Investment/SROI) dan estimasi reduksi jejak karbon ($CO_2$).
%     \end{enumerate}
%     \item Batasan Lingkup Geografis dan Operasional: Pengujian dan implementasi sistem pada fase pilot project dibatasi secara geografis di wilayah Desa Lengkong, Kecamatan Bojongsoang, Kabupaten Bandung. Secara operasional, sistem ini tidak menyediakan armada logistik komersial sendiri. Selain itu, mekanisme pengantaran logistik oleh Relawan kini dibatasi oleh logika batas minimal (threshold); jika makanan yang didonasikan berjumlah sedikit (di bawah 5 porsi), maka opsi pengantaran oleh Relawan akan otomatis dikunci (locked) oleh sistem. Untuk menyiasati donasi dengan kuantitas kecil tersebut, sistem akan mengalihkan pengguna ke metode pengambilan Ambil Sendiri atau fitur Janji Temu (COD) di lokasi netral, yang didukung oleh integrasi jendela pemberitahuan dan tombol WhatsApp untuk memfasilitasi komunikasi langsung antara donatur dan penerima.
%     \item Batasan Teknis dan Keamanan AI: Aplikasi dirancang dengan teknologi Front-End React JS dan Back-End Node/PHP (REST API) menggunakan database MySQL. Aplikasi mewajibkan adanya koneksi internet aktif, akses lokasi (GPS), dan izin kamera keras. Analisis kualitas makanan oleh AI (Food Guard AI) hanya bertindak sebagai bantuan verifikasi visual awal, sedangkan validasi akhir keamanan pangan tetap menjadi tanggung jawab manusia (Admin atau Mitra Penerima) sebelum makanan dikonsumsi.
% \end{enumerate}

\section{Tahapan Pengerjaan}
Dalam pengembangan platform Food AI Rescue, metode rekayasa perangkat lunak yang digunakan adalah pendekatan \textit{Agile} dengan kerangka kerja (\textit{framework}) \textit{Scrum}. Metode ini dipilih karena sifatnya yang iteratif dan fleksibel, sehingga sangat adaptif terhadap perubahan kebutuhan sistem yang dinamis, terutama dalam hal pengintegrasian teknologi kecerdasan buatan (\textit{Artificial Intelligence}) secara \textit{real-time} ke dalam aplikasi. Proses pengembangan dibagi ke dalam siklus-siklus kecil yang berulang untuk memastikan setiap modul dapat dirilis, dievaluasi, dan disempurnakan secara bertahap.

\subsection{Produk \textit{Backlog}}
\textit{Product Backlog} merupakan daftar prioritas utama yang berisi seluruh fitur, fungsi, perbaikan, dan kebutuhan sistem yang harus dibangun dalam platform Food AI Rescue. Daftar ini dikelola secara dinamis untuk memastikan pengembangan terarah pada nilai manfaat aplikasi. Beberapa item utama dalam \textit{Product Backlog} pada proyek ini meliputi:
\begin{enumerate}
    \item Rancang ulang \textit{database}: \textit{Backlog} ini berfokus pada pembaruan struktur penyimpanan data aplikasi dari awal hingga siap beroperasi. Tahapannya sangat komprehensif, dimulai dari merancang desain \textit{database} baru, mengimplementasikannya ke dalam bahasa SQL, hingga melakukan pengujian awal. Setelah kerangka selesai, langkah selanjutnya adalah memasukkan data rekaan (\textit{dummy data}) untuk uji coba, mengoneksikan basis data tersebut ke aplikasi utama, dan melakukan pengujian akhir (\textit{final testing}) untuk memastikan tidak ada kebocoran atau kesalahan aliran data.
    \item Revisi \textit{frontend}: \textit{Backlog} kedua ini bertujuan untuk menyempurnakan antarmuka pengguna (UI) dan pengalaman pengguna (UX) pada aplikasi. Pengerjaannya dimulai dengan menganalisis \textit{bug} atau kecacatan visual yang ada di seluruh modul aplikasi. Setelah masalah teridentifikasi, dilakukan revisi pada kode \textit{frontend} yang dilanjutkan dengan implementasi tata letak dan tampilan baru agar platform Food AI Rescue terlihat lebih rapi, interaktif, dan mudah digunakan oleh para donatur, relawan, maupun penerima.
    \item Implementasi \textit{backend}: \textit{Backlog} ini bergeser ke sisi server untuk memperkuat fondasi logika sistem aplikasi. Tugas utama di sini adalah memigrasikan logika \textit{backend} yang sebelumnya menggunakan \textit{AppScript} ke lingkungan peladen yang baru (seperti Node.js). Selain proses pemindahan, \textit{backlog} ini juga mewajibkan adanya pengujian \textit{backend}, implementasi penanganan eror (\textit{exceptions}) untuk mencegah sistem \textit{crash} saat terjadi anomali, serta pengujian tahap akhir (\textit{final testing}) untuk menjamin stabilitas aplikasi secara keseluruhan.
    \item Implementasi \textit{API Key} AI: \textit{Backlog} terakhir ini merupakan tahap krusial untuk mengaktifkan fitur Kecerdasan Buatan (\textit{Google Gemini API}) yang menjadi inti inovasi Food AI Rescue. Tugasnya meliputi pemasangan atau implementasi \textit{API Key} ke dalam sistem, merevisi konfigurasi \textit{environment} (seperti file \texttt{.env}) untuk memastikan kunci rahasia terlindungi dengan aman, serta melakukan pengujian (\textit{testing} AI) guna memastikan fitur intelijensi visual dan generatif, seperti ``Polisi Pangan'' untuk verifikasi makanan, berjalan dengan akurat.
\end{enumerate}
\subsection{Pelaksanaan \textit{Sprint}}
Pelaksanaan setiap \textit{backlog} dibagi ke dalam beberapa tahapan \textit{sprint} terstruktur guna memastikan pengerjaan yang terukur dan efisien. Rincian tahapan pelaksanaan \textit{sprint} untuk masing-masing \textit{backlog} adalah sebagai berikut:

% \vspace{0.5cm}
\noindent Rancang Ulang \textit{Database}
\begin{itemize}
    \item Desain \textit{database}
    \item Implementasi SQL
    \item \textit{Test database}
    \item \textit{Insert dummy}
    \item Koneksikan dengan aplikasi
    \item \textit{Final testing}
\end{itemize}

\noindent Revisi \textit{Frontend}
\begin{itemize}
    \item Analisis \textit{bug} visual untuk semua modul
    \item Revisi \textit{frontend}
    \item Implementasi tampilan baru
\end{itemize}

\noindent Implementasi \textit{Backend}
\begin{itemize}
    \item Migrasi \textit{logic backend} dari \textit{AppScript}
    \item \textit{Testing backend}
    \item Implementasi \textit{exceptions}
    \item \textit{Final testing}
\end{itemize}

\noindent Implementasi \textit{API Key} AI
\begin{itemize}
    \item Implementasi \textit{API Key}
    \item Revisi \textit{environment}
    \item \textit{Testing} AI
\end{itemize}

\subsection{\textit{Daily Standup}}
Guna menjaga efisiensi komunikasi dan memastikan kelancaran pencapaian target, tim pengembang mengadakan koordinasi rutin melalui mekanisme \textit{Daily Standup}. Rapat evaluasi harian yang singkat ini difokuskan untuk memaparkan perkembangan (\textit{progress}) tugas dari masing-masing personel, menyusun rencana pengembangan kode (\textit{coding}) selanjutnya, serta mendeteksi dan merumuskan solusi bersama terhadap kendala teknis (\textit{blocker}) yang mungkin muncul selama proses integrasi kecerdasan buatan (AI) maupun perancangan arsitektur perangkat lunak.

\subsection{\textit{Scrum Events}}
Untuk menjamin bahwa setiap hasil pengembangan sistem sejalan dengan kebutuhan pengguna, siklus \textit{Sprint} didampingi oleh serangkaian tahapan (\textit{events}) yang terstruktur. Proses ini diawali dengan \textit{Sprint Planning} guna menetapkan prioritas fitur yang akan dibangun. Selanjutnya, \textit{Sprint Review} dilaksanakan pada akhir periode untuk mendemonstrasikan kelayakan fungsi fitur, seperti sistem verifikasi makanan oleh AI dan modul gamifikasi relawan, kepada pemangku kepentingan (\textit{startup} JagoAI). Rangkaian ini kemudian ditutup dengan \textit{Sprint Retrospective} untuk meninjau efektivitas kinerja tim serta merencanakan perbaikan teknis pada iterasi pengembangan berikutnya.

\begin{figure}[H]
    \centering
    \includegraphics[width=1\textwidth]{image/bab1/Scrum.png}
    \caption{Metode Scrum Menggunakan TorMonitorPYanto x Ko+Lab}
    \label{fig:scrum}
\end{figure}

% Pengembangan platform Food AI Rescue (FAR) mengadopsi pendekatan pengerjaan yang iteratif, kolaboratif, dan berbasis \textit{Minimum Viable Product} (MVP). Pelaksanaan program ini dirancang berlangsung selama kurang lebih (waktu pengerjaan) minggu dengan tahapan terstruktur yang berfokus pada adopsi teknologi oleh masyarakat sasaran. Berikut adalah tahapan pengerjaan yang dilakukan:
% \begin{enumerate}
%     \item \textit{Requirement}: \\
%     Tahap ini melibatkan komunikasi dan observasi yang bertujuan untuk memahami perangkat lunak yang dibutuhkan pengguna, termasuk batasan dan ekspektasi platform Food AI Rescue (FAR) sebagai aplikasi non-komersial. Dalam proses ini, informasi dikumpulkan melalui koordinasi dan identifikasi masalah dengan pihak terkait, seperti Mitra Penyedia (UMKM KHB) dan Mitra Penerima (panti asuhan atau komunitas rentan). Informasi tersebut dianalisis untuk merumuskan kebutuhan sistem logistik pangan yang akurat guna memecahkan masalah "keterhubungan yang terputus" (\textit{disconnected link}) secara aman dan efisien.
%     \item \textit{Design}: \\
%     Pada tahap ini, pengembang membuat perancangan sistem yang membantu mendefinisikan arsitektur platform secara keseluruhan. Proses ini mencakup pembuatan desain antarmuka pengguna (\textit{User Interface}) yang mengutamakan aspek keamanan pangan dan kenyamanan navigasi (\textit{ergonomics}), perancangan kerangka \textit{database} relasional menggunakan MySQL, serta pendefinisian logika integrasi dengan layanan pihak ketiga seperti \textit{Application Programming Interface} (API) Gemini AI untuk fitur pemindaian kelayakan visual makanan.
%     \item \textit{Implementation}: \\
%     Pada tahap ini, sistem mulai dikembangkan dan diterjemahkan ke dalam kode program yang berfokus pada pembuatan \textit{Minimum Viable Product} (MVP). Platform Food AI Rescue dikembangkan sebagai \textit{Progressive Web App} (PWA) menggunakan \textit{framework} Node.js. Setiap modul utama, seperti dasbor donatur, sistem titik radius lokasi, keranjang klaim penerima, hingga fitur inti validasi kecerdasan buatan (Food Guard AI), dibangun secara bertahap untuk kemudian diintegrasikan ke dalam satu kesatuan sistem yang utuh.
%     \item \textit{Verification}: \\
%     Pada tahap ini, sistem dilakukan verifikasi dan pengujian operasional \textit{end-to-end} untuk memastikan apakah sistem memenuhi persyaratan fungsionalitas yang diharapkan. Pengujian dilakukan melalui simulasi \textit{Role-Playing} (di mana tim berperan secara langsung sebagai donatur, relawan logistik, dan penerima). Proses ini ditujukan untuk mengevaluasi respons integrasi modul, kelancaran verifikasi otomatis AI, akurasi data \textit{database}, serta mendeteksi adanya kebingungan navigasi (\textit{UX friction}) atau potensi \textit{crash} pada aplikasi.
%     \item \textit{Maintenance}: \\
%     Ini adalah tahap akhir dari metode pengembangan dalam batasan pengerjaan proyek ini. Setelah perangkat lunak diuji coba, pengembang melakukan pemeliharaan berupa stabilisasi sistem akhir. Tahap pemeliharaan ini berfokus pada perbaikan berbagai kesalahan dan \textit{error} minor yang ditemukan pada tahap verifikasi sebelumnya, seperti perbaikan \textit{bug} logika lokasi (peta), penyesuaian tata letak antarmuka, serta penyederhanaan bahasa notifikasi agar aplikasi dinyatakan stabil dan terbebas dari \textit{bug} utama.
% \end{enumerate}


% Halaman ini adalah contoh penggunaan template buku standar LaTeX untuk proyek Book FAR. Struktur template ini dirancang sedemikian rupa agar mudah dikelola dan memenuhi standar estetika modern Telkom University.

% \section{Tujuan Pembuatan Modul}
% Beberapa tujuan utama dari penggunaan template ini adalah:
% \begin{enumerate}
%     \item Meningkatkan efisiensi penulisan buku dan modul praktikum.
%     \item Menciptakan konsistensi styling di seluruh departemen.
%     \item Mempermudah proses pembaruan konten di masa mendatang.
% \end{enumerate}

% \section{Contoh Referensi Gambar}
% Berikut adalah contoh penyisipan gambar yang menggunakan caption terpusat.

% \begin{figure}[ht]
%     \centering
%     % \includegraphics[width=0.5\textwidth]{image/placeholder_figure}
%     \caption{Contoh Gambar dalam Template}
%     \label{fig:contoh_gambar}
% \end{figure}

% Gunakan referensi silang seperti \ref{fig:contoh_gambar} untuk memudahkan pembaca dalam menavigasi konten.

% \section{Contoh Tabel Standar}
% Gunakan paket \texttt{booktabs} untuk membuat tabel yang lebih bersih dan profesional.

% \begin{table}[ht]
%     \centering
%     \caption{Contoh Tabel Data}
%     \label{tab:data_contoh}
%     \begin{tabularx}{\textwidth}{l X X}
%         \toprule
%         No & Komponen & Deskripsi \\
%         \midrule
%         1 & Preamble & Konfigurasi style dan paket utama. \\
%         2 & Frontmatter & Halaman judul dan daftar isi. \\
%         3 & Content & Isi materi per bab. \\
%         \bottomrule
%     \end{tabularx}
% \end{table}

% \section{Kesimpulan}
% Dengan menggunakan template modular ini, Anda dapat fokus pada isi konten tanpa harus sering-sering mengedit preamble untuk setiap bab.
